\section{What Zhu proves}\label{zh:sec:zhu}

Zhu \cite{Zhu26} writes $L$ for the half length of the window. Here it is $a$, as in Chapter~\ref{chap:yoshida}, where $L$ is the support of an autocorrelation, so that his $\lambda^*(L)$, $A_L$ and $T_1(L)$ appear below as $\lambda^{\rm even}_{\min}(\frac12,a)$, $A_a$ and $T_1(a)$. For real even continuous $f$ supported in $[-a,a]$, with $F=\widehat f$, $\widehat f(t)=\int_\R f(y)e^{ity}\,dy$, Zhu writes Weil's quadratic form as
\begin{equation}\label{zh:Q}
Q(f)=2F\bigl(\tfrac i2\bigr)^2+\frac1{2\pi}\int_\R|F(t)|^2\,\Psi_a(t)\,dt,\qquad
\Psi_a(t)=\re\psi\Bigl(\frac14+\frac{it}2\Bigr)-\log\pi-\mathcal P_a(t),
\end{equation}
where $\mathcal P_a(t)=\sum_{\log n<2a}2\Lambda(n)n^{-1/2}\cos(t\log n)$ is the prime comb. He sets $\lambda^{\rm even}_{\min}(\frac12,a)=\inf Q(f)/\|f\|^2$ over the real even $f\ne0$ supported in $[-a,a]$. His results, with the status he gives them, are the following.
\begin{itemize}
\item[(Z1)] One stroke reduction (Theorem 1.1, unconditional). Let $A_a=\sup_t\mathcal P_a(t)=\sum_{\log n<2a}2\Lambda(n)n^{-1/2}$: the supremum is attained at $t=0$, and by Weyl equidistribution of the phases $t\log p$ values arbitrarily close to it occur at arbitrarily large $t$ (his Lemma 3.2). The envelope $\re\psi(\frac14+\frac{it}2)-\log\pi\ge\log\frac t{2\pi}-\frac1t$ for $t\ge\frac{15}4$ (Lemma 3.1) gives $\Psi_a(t)\ge\beta^*>0$ for $t\ge T_\sharp$, where $\beta^*=\log\frac{T_\sharp}{2\pi}-\frac1{T_\sharp}-A_a$ and $\beta^*>0$ holds for every $T_\sharp$ above a threshold that exceeds $T_1:=2\pi e^{A_a}$ by the factor $1+O(T_1^{-1})$ ($120.1$ against $119.1$ at $a=0.8$). Positivity on the window then reduces to the positive semidefiniteness of a finite matrix in an even Legendre basis, with explicit tails.
\item[(Z2)] Certificate (Theorem 1.2, a computer assisted theorem). At $a=0.8$, with $T_\sharp=200$, $200$ even Legendre modes (orders up to $398$), working precision $50$ digits and a verified Cholesky factorisation, $Q(f)\ge8.9\times10^{-18}\|f\|^2$ for every real even $f$ supported in $[-0.8,0.8]$. The autocorrelation then has support $[-1.6,1.6]$, which is $2.3$ times the range $2a\le\log2$ that Zhu attributes to Yoshida \cite{Yoshida92} and, by a different method, to Connes and Consani \cite{ConnesConsani21}. His variational upper bound, computed on the geometric side in interval arithmetic, is $\lambda^{\rm even}_{\min}(\frac12,0.8)\le2.27\times10^{-17}$, which encloses the infimum within a factor $2.6$.
\item[(Z3)] Under RH, $\lambda^{\rm even}_{\min}(\frac12,a)\le\exp(-ae^a)$ for $a$ large (Theorem 1.3, conditional on RH), by an explicit test function: a canonical product over the ordinates up to $T=2\pi e^a$ times a power of the sinc function. The construction works at half saturation, and with $T=2\pi e^{2\theta a}$ it closes only for $\theta\lesssim0.65$ (his Remark 13.6).
\item[(Z4)] Decay law (Conjecture 12.1, fitted to computed bounds). The certified variational upper bounds for $a$ from $0.5$ to $2$, unconditional and computed on the geometric side, fall from $10^{-6}$ to $3.2\times10^{-283}$ and follow $-\log\lambda^{\rm even}_{\min}(\frac12,a)=C\,N(T_*)/\log N(T_*)\,(1+o(1))$, where $T_*=2\pi e^{2a}$, $N(T)$ is the number of ordinates up to $T$ and $C=20.13\ldots\approx2\pi^2$ (the fit of his equation (5)). The constant is fitted and not derived: Zhu calibrates the law on $a\le1.2$ and tests it on $a=1.4$ to $2$, where it underpredicts by about $5$ percent with no trend. That the law describes the infimum itself, with $2\pi^2$ in place of the fitted $C$, is his Conjecture 12.1.
\item[(Z5)] Barrier (Theorem 1.4, unconditional). Since $A_a=(4+o(1))e^a$ by the prime number theorem, any application of (Z1) needs $T_\sharp>T_1(a)=2\pi e^{A_a}$, about $1.0\times10^7$ at $a=1.6$ and $2.4\times10^{11}$ at $a=2$ ($A_{1.6}=14.323$, $A_2=24.383$), so that the size of the matrix grows doubly exponentially in $a$; and since $\sup_t\mathcal P_a(t)=A_a$ exactly, no pointwise bound on the comb lowers $T_1$. To improve on pointwise envelopes one would have to prove that the phases $t\log p$ cannot align, which Zhu calls an arithmetic statement about simultaneous Diophantine approximation of $\{\log p\}$ (his Remark 1.6).
\item[(Z6)] Parity (Lemma 6.1, Theorem 6.2 and Corollary 6.3). For real $f=f_e+f_o$, $Q(f)=Q(f_e)+Q(f_o)$, with the pole term entering as $+2c^2$ on the even part and $-2s^2$ on the odd part, where $c=\int f_e(x)\cosh\frac x2\,dx$ and $s=\int f_o(x)\sinh\frac x2\,dx$. At $a=0.8$ the even ground state is simple, with the certified bounds $8.9\times10^{-18}\le\lambda_1^{\rm even}\le2.523\times10^{-16}$ (the upper bound $2.27\times10^{-17}$ of (Z2) is the sharper), $\lambda_2^{\rm even}\ge2.085\times10^{-12}$ and $8.206\times10^{-15}\le\lambda_1^{\rm odd}\le2.347\times10^{-14}$, where $\lambda_1^{\rm even}\le\lambda_2^{\rm even}$ are the first two min max values of $Q(f)/\|f\|^2$ over real even $f$ supported in $[-0.8,0.8]$ and $\lambda_1^{\rm odd}$ is the infimum over real odd $f$. Hence $Q(f)\ge8.9\times10^{-18}\|f\|^2$ for complex $f$ as well. Zhu connects the simple even ground state with one of the two spectral steps that the programme of Connes, Consani and Moscovici \cite{CCM25} isolates as missing.
\end{itemize}
Zhu also retracts a claim of an earlier draft at autocorrelation support $2.38$ ($a=1.19$), keeping a computation that is consistent with positivity there and no certificate (his Section 7), and he warns that evaluating the geometric side of $Q$ in double precision produces spurious negative eigenvalues (his Remark 2.1 and Section 10).

For complex test functions the least eigenvalue of Chapter~\ref{chap:yoshida} is $\lambda_{\min}(\frac12,a)$, the infimum of $\langle\varphi,\varphi\rangle/\|\varphi\|^2$ over $0\ne\varphi\in C(a)$ (Section~\ref{yo:sec:coercive}). Write $\lambda^{\rm odd}_{\min}(\frac12,a)$ for the infimum over real odd test functions. The form has no cross terms between the real and the imaginary part of $\varphi$, nor between its even and its odd part (Zhu's Lemma 6.1), so that
\[
\lambda_{\min}(\tfrac12,a)=\min\bigl(\lambda^{\rm even}_{\min}(\tfrac12,a),\,\lambda^{\rm odd}_{\min}(\tfrac12,a)\bigr),
\]
and the two sides $\lambda_{\min}(\frac12,a)$ and $\lambda^{\rm even}_{\min}(\frac12,a)$ are equal when the ground state is even, which Zhu's Theorem 6.2 shows at $a=0.8$. Yoshida's parity reduction (Y1) of Section~\ref{yo:sec:yoshida} concerns the same two sectors. The splitting holds for the damped forms of Section~\ref{yo:sec:mult} at every $\alpha$ as well, because the damping $e^{-\varepsilon|y|}$ is even.

\section{The Thorin dictionary}\label{zh:sec:dict}

The functional $\mathcal Y$, the form $\langle\varphi_1,\varphi_2\rangle=\mathcal Y(\varphi_1*\widetilde\varphi_2)$ and the spaces $C(a)$ and $K(a)$ are those of Section~\ref{yo:sec:yoshida}. Call a real even $f$ supported in $[-a,a]$ admissible if the explicit formula applies to $f*f$, as it does for every $f\in C_c^\infty(\R)$. Let $PW_a$ be the Paley Wiener space of the entire functions of exponential type at most $a$ that are square integrable on $\R$, and $PW_a^{\rm even}$ its real even members, which contain the transforms $F=\widehat f$ of the admissible $f$. The measure $U=\sum_{\gamma>0}\delta_{\gamma^2}$ is the Thorin measure of the clock under RH (Section~\ref{ch:clock}), each ordinate counted with its multiplicity.

\begin{proposition}\label{zh:prop:dict}
For admissible $f$, with $F=\widehat f$, Zhu's form is Yoshida's form on the window $C(a)$: $Q(f)=\langle f,f\rangle=\mathcal Y(f*f)$. Under RH
\[
\begin{gathered}
Q(f)=\sum_\rho|F(\gamma_\rho)|^2=2\int_{(0,\infty)}|F(\sqrt t)|^2\,U(dt),\\
\lambda^{\rm even}_{\min}(\tfrac12,a)=\inf_f\frac{2\int_{(0,\infty)}|F(\sqrt t)|^2\,U(dt)}{\frac1{2\pi}\int_\R|F|^2},
\end{gathered}
\]
so that $\lambda^{\rm even}_{\min}(\frac12,a)$ is the lower sampling constant of the Thorin measure $U$ for the transforms $F\in PW_a^{\rm even}$ of the admissible $f$.
\end{proposition}

\begin{proof}
The autocorrelation $g=f*f$ is even, supported in $[-2a,2a]$ and has $\widehat g=F^2=|F|^2\ge0$ on $\R$. By Proposition~\ref{yo:prop:Z} (there $L=2a$; the proposition is stated for smooth $f$, and for an admissible $f$ the identity is the term by term explicit formula that follows), with no hypothesis on the zeros, $\langle f,f\rangle=\mathcal Y(g)=2\int\widehat g(\sqrt t)\,U^\pm(dt)$, and only the prime powers $n\le e^{2a}$ enter; under RH the functional $\varphi\mapsto\int\varphi\,dU^\pm$ is the sum of $\varphi$ over the atoms of $U$ (Section~\ref{sec:ch14:functional}). Term by term, the explicit formula \eqref{yo:EF} for $g$ has the pole term $\widehat g(\frac i2)+\widehat g(-\frac i2)=2F(\frac i2)^2$. Its prime term $-\sum_{n\ge2}\Lambda(n)n^{-1/2}\bigl(g(\log n)+g(-\log n)\bigr)$ equals $-\frac1{2\pi}\int|F|^2\mathcal P_a$, because $\widehat g=|F|^2$ and $g$ vanishes beyond $2a$. Its archimedean term is $V_1(g)-g(0)\log\pi=\frac1{2\pi}\int|F|^2\bigl(\re\psi(\frac14+\frac{it}2)-\log\pi\bigr)dt$, by the first line of \eqref{yo:V1} and $g(0)=\frac1{2\pi}\int|F|^2$. The three terms add up to \eqref{zh:Q}. Under RH every $\gamma_\rho$ is real, and the zero side of \eqref{yo:EF} is $\sum_\rho\widehat g(\gamma_\rho)=\sum_\rho F(\gamma_\rho)^2=2\sum_{\gamma>0}F(\gamma)^2$, because $F$ is even and the ordinates come in pairs $\pm\gamma$; this is $2\int F(\sqrt t)^2\,U(dt)$. Finally $\|f\|^2=\frac1{2\pi}\int|F|^2$ by Plancherel.
\end{proof}

Two consequences follow. First, under RH the positivity of the infimum for every $a$ is a sampling statement, and it is Zhu's Proposition 2.3. Littlewood's theorem on the gaps between consecutive ordinates \cite[Section~9.12]{Titchmarsh}, which holds without RH although Zhu states it as a consequence of RH, says that the gaps tend to zero, so under RH the set $\{\pm\gamma\}$ contains, above some height, a separated subset $\Lambda$ of lower uniform density larger than $a/\pi$ (a maximal subset of points at least $\pi/(4a)$ apart), and Beurling's sampling theorem \cite[Chapters 2 and 3]{Seip04} gives $c>0$ with $\sum_{\lambda\in\Lambda}|F(\lambda)|^2\ge c\int_\R|F|^2$ for all $F\in PW_a$. Hence $Q(f)\ge2\pi c\|f\|^2$ and $\lambda^{\rm even}_{\min}(\frac12,a)>0$, with a constant that is not explicit (Zhu's Remark 2.4). Without RH the argument gives nothing (Zhu's Remark 2.4 counts the gap bound as a second use of RH and says that unconditionally the gaps are only known to be bounded). Second, in this picture the smallness of the infimum is the poor conditioning of sampling by a set that is sparse at low frequency: there is no ordinate in $(-14.13,14.13)$.
% CHECK: Littlewood's theorem on the gaps between consecutive ordinates is unconditional, gamma_{n+1} - gamma_n = O(1/log log log gamma_n) (Littlewood 1924; so stated in the introduction of arXiv 1410.3635 and in OEIS A208436), and the improvement O(1/log log gamma_n) is the one that needs RH (Titchmarsh Chapter XIV, recalled); Zhu calls the theorem a consequence of RH, cites Ch. XIII of Titchmarsh and says in his Remark 2.4 that unconditionally the gaps are only known to be bounded (finding M4 of misc/reviews/nick-zhu/final-Z18.md, settled against the two sources named). The pinpoints Titchmarsh Section 9.12 (taken from review-A.md Section 6 with medium confidence) and Seip Chapters 2 and 3 (Zhu's pinpoint for Beurling's theorem) are not checked against the books; the bibliography entry Seip04 is in misc/reviews/nick-zhu/bib-new.tex.

Zhu's parity lemma (Z6) has a Thorin reading. The pole is the atom of mass one of the signed Thorin measure $U^\pm$ at $t=-\frac14$ (equation~\eqref{yo:signed}), that is at the imaginary ordinates $\pm\frac i2$, and it is paired with $F(\frac i2)F(-\frac i2)$. For even $f$ this is $c^2\ge0$ and for odd $f$ it is $-s^2\le0$ (Proposition~\ref{yo:prop:Z}): it helps the even sector and hurts the odd one. That the even ground state is nevertheless the smaller is explained by Zhu (Section 6.3) through the constraint $F_o(0)=0$ on the odd sector, which denies odd test functions the zero free interval $[0,\gamma_1)$ in which even minimisers park most of their mass: the low frequency gap, and not the pole, controls the conditioning. At $a=0.8$ the odd sector lies about three orders of magnitude above the even one and at $a=0.5$ two orders (Table~\ref{zh:tab:hp}), and Zhu reports that the ratio grows almost geometrically in $a$ over $[0.5,1.2]$, as $10^{2.07a+1.32}$.

The least eigenvalues of the two sectors are computed from the prime side matrices of Yoshida's form. The even sector has the basis $(2a)^{-1/2}$ and $a^{-1/2}\cos(\pi jy/a)$, the odd sector the basis $a^{-1/2}\sin(\pi jy/a)$, for $1\le j\le N$, and $N$ is the number of modes in Table~\ref{zh:tab:hp}.

\begin{table}[t]
\centering
\fitwidth{\begin{tabular}{rrrrr}
\toprule
$a$ & modes $N$ & $\lambda^{\rm even}_{\min}$ & $\lambda^{\rm odd}_{\min}$ & Zhu, even\\
\midrule
0.5 & 8 & $1.336\times10^{-6}$ & $2.51\times10^{-4}$ & \\
0.5 & 14 & $1.082\times10^{-6}$ & $2.39\times10^{-4}$ & \\
0.5 & 20 & $1.005\times10^{-6}$ & $2.30\times10^{-4}$ & \\
0.5 & 26 & $9.743\times10^{-7}$ & $2.24\times10^{-4}$ & $\le1.05\times10^{-6}$\\
0.5 & 160 & $9.344\times10^{-7}$ & $1.997\times10^{-4}$ & \\
\midrule
0.8 & 8 & $5.49\times10^{-15}$ & $2.67\times10^{-12}$ & \\
0.8 & 14 & $4.127\times10^{-17}$ & $3.86\times10^{-14}$ & \\
0.8 & 20 & $2.651\times10^{-17}$ & $2.333\times10^{-14}$ & \\
0.8 & 26 & $2.144\times10^{-17}$ & $2.157\times10^{-14}$ & $[8.9\times10^{-18},\,2.27\times10^{-17}]$\\
0.8 & 160 & $1.673\times10^{-17}$ & $1.644\times10^{-14}$ & \\
\bottomrule
\end{tabular}}
\caption{Least eigenvalues of Weil's form on the window $[-a,a]$ in the even and the odd sector, from the prime side Yoshida matrices in $60$ digit arithmetic with cosine and sine modes up to the order $N$ shown (\texttt{code/zh\_window\_eigs.py}). The last column gives Zhu's certified bounds for the even sector. The values at $26$ modes and below are not converged, and the odd values at $160$ modes are still falling.}
\label{zh:tab:hp}
\end{table}
% CHECK: Zhu's Table 1 reference floors at three digits are not reproduced by the finite sections of data/zh_window_eigs.json, up to 400 modes, in three entries: a = 0.5 odd (Zhu 1.94e-4; ours 1.962e-4 at 400 modes, extrapolation 1.937e-4), a = 0.8 even (Zhu 1.65e-17; ours 1.640e-17, extrapolation 1.6345e-17, below his value) and a = 0.8 odd (Zhu 1.57e-14; ours 1.596e-14, still falling); the even value at a = 0.5 (Zhu 9.34e-7; ours 9.335e-7) and the second even value at a = 0.8 (Zhu 8.38e-12; ours 8.301e-12) also lie below his. The paragraph after Table zh:tab:hp compares at 160 modes only, and whether it should say what the sections of 192 to 400 modes show is for the authors (misc/reviews/nick-zhu/integration-Z3.md, Section 5, points 1 and 2).

The values of Table~\ref{zh:tab:hp} are consistent with Zhu's numbers, which come from another basis and another code: Fourier modes on $[-a,a]$ against his Legendre polynomials, and quadrature of the explicit formula in the variable $y$ against his spherical Bessel integrals. At $a=0.8$ both sectors lie inside his certified intervals at $26$ modes. At $160$ modes the even values $9.344\times10^{-7}$ and $1.673\times10^{-17}$ lie $0.04$ and $1.4$ percent above the reference floors $9.34\times10^{-7}$ and $1.65\times10^{-17}$ of Zhu's Table 1 (a converged exact assembly, which he gives as reference values and not as certificates), and the odd values $1.997\times10^{-4}$ and $1.644\times10^{-14}$ lie $2.9$ and $4.7$ percent above his $1.94\times10^{-4}$ and $1.57\times10^{-14}$ and are still falling. At $a=0.8$ the second even eigenvalue is $8.42\times10^{-12}$ at $160$ modes (Zhu: $8.38\times10^{-12}$ in his Table 1, and the certified bound $\lambda_2^{\rm even}\ge2.085\times10^{-12}$), about $5\times10^5$ times the first, in agreement with the simple ground state of (Z6). Raising the precision to $90$ digits, or the number of quadrature nodes by half, changes the least eigenvalues by a relative $2\times10^{-58}$ at most (checked at $26$ and $40$ modes), so every value is an upper bound for the infimum over its sector.

As a check of the normalisation of Proposition~\ref{zh:prop:dict}, take the unit bottom eigenvector $f$ of the matrices of $26$ and $40$ modes at $a=0.5$ and $0.8$, in both sectors. The partial sums $2\sum_{\gamma\le\gamma_K}F(\gamma)^2$ over the first $K$ ordinates of \texttt{data/zeros\_1700.txt} increase with $K$ and stay below the eigenvalue, and with the mean tail $2\int_{\gamma_K}^\infty F(t)^2\,\frac1{2\pi}\log\frac t{2\pi}\,dt$ of the missing ordinates they reproduce it to within $0.16$ percent for $K=300$, $400$ and $1700$ in all eight cases.

Three other tables of the book list the same eigenvalue at other levels: Table~\ref{yo:tab:lam} of Chapter~\ref{chap:yoshida} for the windows $a=0.35$ to $4$ (its entry $2.28\times10^{-17}$ at $a=0.8$, from $24$ modes, lies between our values at $20$ and $26$ modes), and Tables~\ref{dc:tab:ladder} and~\ref{dc:tab:zero} of Chapter~\ref{chap:dc}. Table~\ref{zh:tab:hp} is at $\alpha=\frac12$ only, with the even and the odd sector apart, in $60$ digits, and it shows the convergence in the number of modes without extrapolating; it reaches $a=0.8$.
